\begin{table}[H]
\centering
\small
\caption{Hiperparâmetros: grades e escolhas da validação cruzada embargada.}
\label{tab:prev-hiperparametros}
\begin{tabular}{lp{6.2cm}lc}
\toprule
Modelo & Grade & Mais frequente & No limite \\
\midrule
Ridge & $\alpha$: 17 valores de $10^{-3}$ a $10^{5}$ & $\alpha = 10^{3}$ (24) & 0 \\
LASSO & $\alpha$: 16 valores de $10^{-3}$ a $10^{2}$ & $\alpha = 4{,}64$ (18) & 0 \\
\addlinespace
Floresta aleatória & profundidade máxima: 1, 2, 3, 6, sem limite & 1 (33) & 33 \\
 & folha mínima: 5, 20, 50, 100, 200 & & \\
 & fração de variáveis por corte: 0{,}2; 1/3; 1 & & \\
\addlinespace
XGBoost & taxa de aprendizado: 0{,}01; 0{,}03; 0{,}1 & 0{,}01 (33) & 33 \\
 & profundidade máxima: 1, 2, 4 & 1 (28) & 28 \\
 & número de árvores: 50, 100, 200, 500 & 50 (22) & 22 \\
 & peso mínimo por folha: 1, 10 & & \\
\bottomrule
\end{tabular}
\par\smallskip
\parbox{0.95\linewidth}{\footnotesize\textit{Nota}: Janela expansiva, 33 reestimações. Em cada uma, validação cruzada com cinco dobras expansivas e embargo de 30 dias; escolhe-se a combinação de menor erro quadrático médio. Entre parênteses, o número de reestimações com a escolha indicada. Última coluna: reestimações com o hiperparâmetro no extremo de menor capacidade da grade. Floresta com 500 árvores; XGBoost com 80\% das observações e das variáveis em cada árvore. Os modelos lineares usam variáveis padronizadas com a média e o desvio-padrão do treino.}
\end{table}
